% Options for packages loaded elsewhere
\PassOptionsToPackage{unicode}{hyperref}
\PassOptionsToPackage{hyphens}{url}
\PassOptionsToPackage{dvipsnames,svgnames,x11names}{xcolor}
\documentclass[
  10pt,
  fontset=fandol]{ctexart}
\usepackage{xcolor}
\usepackage[margin=16mm]{geometry}
\usepackage{amsmath,amssymb}
\setcounter{secnumdepth}{-\maxdimen} % remove section numbering
\usepackage{iftex}
\ifPDFTeX
  \usepackage[T1]{fontenc}
  \usepackage[utf8]{inputenc}
  \usepackage{textcomp} % provide euro and other symbols
\else % if luatex or xetex
  \usepackage{unicode-math} % this also loads fontspec
  \defaultfontfeatures{Scale=MatchLowercase}
  \defaultfontfeatures[\rmfamily]{Ligatures=TeX,Scale=1}
\fi
\usepackage{lmodern}
\ifPDFTeX\else
  % xetex/luatex font selection
\fi
% Use upquote if available, for straight quotes in verbatim environments
\IfFileExists{upquote.sty}{\usepackage{upquote}}{}
\IfFileExists{microtype.sty}{% use microtype if available
  \usepackage[]{microtype}
  \UseMicrotypeSet[protrusion]{basicmath} % disable protrusion for tt fonts
}{}
\makeatletter
\@ifundefined{KOMAClassName}{% if non-KOMA class
  \IfFileExists{parskip.sty}{%
    \usepackage{parskip}
  }{% else
    \setlength{\parindent}{0pt}
    \setlength{\parskip}{6pt plus 2pt minus 1pt}}
}{% if KOMA class
  \KOMAoptions{parskip=half}}
\makeatother
\usepackage{graphicx}
\makeatletter
\newsavebox\pandoc@box
\newcommand*\pandocbounded[1]{% scales image to fit in text height/width
  \sbox\pandoc@box{#1}%
  \Gscale@div\@tempa{\textheight}{\dimexpr\ht\pandoc@box+\dp\pandoc@box\relax}%
  \Gscale@div\@tempb{\linewidth}{\wd\pandoc@box}%
  \ifdim\@tempb\p@<\@tempa\p@\let\@tempa\@tempb\fi% select the smaller of both
  \ifdim\@tempa\p@<\p@\scalebox{\@tempa}{\usebox\pandoc@box}%
  \else\usebox{\pandoc@box}%
  \fi%
}
% Set default figure placement to htbp
\def\fps@figure{htbp}
\makeatother
\setlength{\emergencystretch}{3em} % prevent overfull lines
\providecommand{\tightlist}{%
  \setlength{\itemsep}{0pt}\setlength{\parskip}{0pt}}
\usepackage{amsmath,amssymb,mathtools,needspace}
\xeCJKDeclareCharClass{CJK}{"2032,"0400->"04FF,"2160->"216B,"2460->"2473,"25A0,"25C7,"25CB,"25CF}
\IfFontExistsTF{Arial Unicode MS}{\xeCJKsetup{AutoFallBack=true}\setCJKfallbackfamilyfont{\CJKrmdefault}{Arial Unicode MS}}{}
\setcounter{secnumdepth}{0}
\setlength{\emergencystretch}{2em}
\allowdisplaybreaks
\usepackage{bookmark}
\IfFileExists{xurl.sty}{\usepackage{xurl}}{} % add URL line breaks if available
\urlstyle{same}
\hypersetup{
  pdftitle={线性插值与抛物插值},
  colorlinks=true,
  linkcolor={Maroon},
  filecolor={Maroon},
  citecolor={Blue},
  urlcolor={Blue},
  pdfcreator={LaTeX via pandoc}}

\title{线性插值与抛物插值}
\author{}
\date{}

\begin{document}
\maketitle

\subsubsection{2.2.2 线性插值与抛物插值}\label{ux7ebfux6027ux63d2ux503cux4e0eux629bux7269ux63d2ux503c}

从定理 2.1 的证明可看到，要求插值多项式 \(P(x)\)，可以通过求方程组（2.2.1）的解 \(a_0,a_1,\cdots,a_n\) 得到。但这样做不但计算复杂，且难以得到 \(P(x)\) 的简单表达式。为了求得便于使用的简单的插值多项式 \(P(x)\)，先讨论 \(n=1\) 的情形。

假定已知区间 \([x_k,x_{k+1}]\) 的端点处的函数值 \(y_k=f(x_k)\)，\(y_{k+1}=f(x_{k+1})\)，要求线性插值多项式 \(L_1(x)\)，使它满足

\[
L_1(x_k)=y_k,\qquad L_1(x_{k+1})=y_{k+1}.
\]

\(y=L_1(x)\) 的几何意义就是通过两点 \((x_k,y_k)\) 和 \((x_{k+1},y_{k+1})\) 的直线，如图 2.2 所示。\(L_1(x)\) 的表达式可由几何意义直接给出，即

\[
L_1(x)=y_k+\frac{y_{k+1}-y_k}{x_{k+1}-x_k}(x-x_k)
\quad\text{（点斜式）},
\tag{2.2.3}
\]

\[
L_1(x)=\frac{x_{k+1}-x}{x_{k+1}-x_k}y_k
+\frac{x-x_k}{x_{k+1}-x_k}y_{k+1}
\quad\text{（两点式）}.
\tag{2.2.4}
\]

\includegraphics[width=0.85\linewidth,height=\textheight,keepaspectratio,alt={图 2.2（原图裁切）}]{../assets/ch02a/figure-2.2.png}

图 2.2（原图）：直线 \(y=L_1(x)\) 与曲线 \(y=f(x)\) 在两节点相交。全部文字标签：\(y\)、\(x\)、\(O\)、\(x_k\)、\(x_{k+1}\)、\(y_k\)、\(y_{k+1}\)、\(y=L_1(x)\)、\(y=f(x)\)。

由两点式方程看出，\(L_1(x)\) 是由两个线性函数

\[
l_k(x)=\frac{x-x_{k+1}}{x_k-x_{k+1}},
\qquad
l_{k+1}(x)=\frac{x-x_k}{x_{k+1}-x_k}
\]

的线性组合得到的，其系数分别为 \(y_k\) 及 \(y_{k+1}\)，即

\[
L_1(x)=y_kl_k(x)+y_{k+1}l_{k+1}(x).
\tag{2.2.5}
\]

显然，\(l_k(x)\) 及 \(l_{k+1}(x)\) 也是线性插值多项式，在节点 \(x_k\) 及 \(x_{k+1}\) 上满足条件

\[
l_k(x_k)=1,\quad l_k(x_{k+1})=0,\quad
l_{k+1}(x_k)=0,\quad l_{k+1}(x_{k+1})=1.
\]

称函数 \(l_k(x)\) 及 \(l_{k+1}(x)\) 为\textbf{一次插值基函数}或\textbf{线性插值基函数}，它们的图形分别如图 2.3(a)、(b) 所示。

\includegraphics[width=0.85\linewidth,height=\textheight,keepaspectratio,alt={图 2.3（原图裁切）}]{../assets/ch02a/figure-2.3.png}

图 2.3（原图）：(a) \(l_k(x)\) 在 \(x_k\) 处为 \(1\)，在 \(x_{k+1}\) 处为 \(0\)；(b) \(l_{k+1}(x)\) 在 \(x_k\) 处为 \(0\)，在 \(x_{k+1}\) 处为 \(1\)。两图全部文字标签：\(y\)、\(x\)、\(O\)、\(1\)、\(x_k\)、\(x_{k+1}\)、\(l_k(x)\)、\(l_{k+1}(x)\)、(a)、(b)。

下面讨论 \(n=2\) 的情况。假定插值节点为 \(x_{k-1},x_k,x_{k+1}\)，要求二次插值多项式 \(L_2(x)\)，使它满足

\begin{quote}
页末插值条件续 PDF 30。
\end{quote}

\href{../../source-images/pdf-030.jpeg}{原页：书页 17／PDF 30}

\begin{quote}
页首承接 PDF 29 的二次插值条件；末式展开续 PDF 31。
\end{quote}

\[
L_2(x_j)=y_j\qquad(j=k-1,k,k+1).
\]

我们知道，\(y=L_2(x)\) 在几何上就是通过三点 \((x_{k-1},y_{k-1})\)、\((x_k,y_k)\)、\((x_{k+1},y_{k+1})\) 的抛物线。为了求出 \(L_2(x)\) 的表达式，可采用基函数方法，此时基函数 \(l_{k-1}(x),l_k(x)\) 及 \(l_{k+1}(x)\) 是二次函数，且在节点上满足条件

\[
\begin{cases}
l_{k-1}(x_{k-1})=1,\quad l_{k-1}(x_j)=0\quad(j=k,k+1),\\
l_k(x_k)=1,\quad l_k(x_j)=0\quad(j=k-1,k+1),\\
l_{k+1}(x_{k+1})=1,\quad l_{k+1}(x_j)=0\quad(j=k-1,k).
\end{cases}
\tag{2.2.6}
\]

满足条件式（2.2.6）的插值基函数是很容易求出的，例如求 \(l_{k-1}(x)\)，因它有两个零点 \(x_k\) 及 \(x_{k+1}\)，故可表示为

\[
l_{k-1}(x)=A(x-x_k)(x-x_{k+1}),
\]

其中 \(A\) 为待定系数，可由条件 \(l_{k-1}(x_{k-1})=1\) 定出，即

\[
A=\frac{1}{(x_{k-1}-x_k)(x_{k-1}-x_{k+1})},
\]

故

\[
l_{k-1}(x)=
\frac{(x-x_k)(x-x_{k+1})}
{(x_{k-1}-x_k)(x_{k-1}-x_{k+1})}.
\]

同理，

\[
l_k(x)=
\frac{(x-x_{k-1})(x-x_{k+1})}
{(x_k-x_{k-1})(x_k-x_{k+1})},
\qquad
l_{k+1}(x)=
\frac{(x-x_{k-1})(x-x_k)}
{(x_{k+1}-x_{k-1})(x_{k+1}-x_k)}.
\]

函数 \(l_{k-1}(x),l_k(x),l_{k+1}(x)\) 称为\textbf{二次插值基函数}或\textbf{抛物插值基函数}，它们在区间 \([x_{k-1},x_{k+1}]\) 上的图形分别如图 2.4(a)、(b)、(c) 所示。

\includegraphics[width=0.85\linewidth,height=\textheight,keepaspectratio,alt={图 2.4（原图裁切）}]{../assets/ch02a/figure-2.4.png}

图 2.4（原图）：(a) \(l_{k-1}(x)\)，(b) \(l_k(x)\)，(c) \(l_{k+1}(x)\)，分别在对应节点取 \(1\)、其他两个节点取 \(0\)；左右基函数在部分区间为负。三图全部文字标签：\(y\)、\(x\)、\(O\)、\(1\)、\(x_{k-1}\)、\(x_k\)、\(x_{k+1}\)、\(l_{k-1}(x)\)、\(l_k(x)\)、\(l_{k+1}(x)\)、(a)、(b)、(c)。

利用二次插值基函数 \(l_{k-1}(x),l_k(x),l_{k+1}(x)\)，立即可得到二次插值多项式

\[
L_2(x)=y_{k-1}l_{k-1}(x)+y_kl_k(x)+y_{k+1}l_{k+1}(x),
\tag{2.2.7}
\]

显然，它满足条件

\[
L_2(x_j)=y_j\qquad(j=k-1,k,k+1).
\]

将上面求得的 \(l_{k-1}(x),l_k(x),l_{k+1}(x)\) 代入式（2.2.7），得

\href{../../source-images/pdf-031.jpeg}{原页：书页 18／PDF 31}

\begin{quote}
页首续 PDF 30 的式（2.2.7）展开；页末公式续 PDF 32。
\end{quote}

\[
\begin{aligned}
L_2(x)={}&y_{k-1}
\frac{(x-x_k)(x-x_{k+1})}
{(x_{k-1}-x_k)(x_{k-1}-x_{k+1})}\\
&+y_k
\frac{(x-x_{k-1})(x-x_{k+1})}
{(x_k-x_{k-1})(x_k-x_{k+1})}\\
&+y_{k+1}
\frac{(x-x_{k-1})(x-x_k)}
{(x_{k+1}-x_{k-1})(x_{k+1}-x_k)}.
\end{aligned}
\]

\end{document}
